\section{Introduction}

Resource estimation often sits between a fast analytical formula and an expensive optimizer. A formula is attractive because it can be evaluated before compilation, but its authority is easy to overstate. Zero error on a benchmark is not a theorem. Conversely, a theorem about a broad exact family does not automatically provide a compact explanation of which structural regime produced a cost.

The TARE regime programme offers a useful separation. An initial three-family expression had a clear semantic interpretation: donor-exact, split-anchor, or central borrow. It matched exhaustive and held-out panels and was used for prospective staging. QG-5 then did what a useful certification system should allow: it found one exact fresh counterexample and recorded the forecast as wrong rather than weakening the gate.

The later support-two theorem makes a stronger repair possible. Instead of adding a case to the old formula and calling the fit complete, we evaluate the exact support-$\le2$ family whose equality with unrestricted DP is already proved for all $n$ under the frozen objective. This produces a theorem-backed exact cost forecaster $F_2$. Compact regime explanations can then be developed and refuted independently on top of that cost layer.

This paper argues for that layered view of a certified forecast. We distinguish the cost theorem, the implementation check against DP, the human-readable regime model, prospective prediction evidence, and row-level verification authority. The QG-7 refutation of the first repaired explanation is valuable precisely because the layered contract prevents it from being mistaken for a failure of the exact cost certificate.